\setcounter{section}{3}








Das
\definitionswort {Kroneckerprodukt}{}
zu Matrizen

\mavergleichskettedisp
{\vergleichskette
{ A
}
{ =} { \left(  a_{ij}  \right)_{1 \leq i \leq m,\, 1 \leq j \leq n}
}
{ } {
}
{ } {
}
{ } {
}
}
{}{}{}
und

\mavergleichskettedisp
{\vergleichskette
{ B
}
{ =} { \left(  b_{k \ell }  \right)_{1 \leq k \leq p,\, 1 \leq \ell \leq r}
}
{ } {
}
{ } {
}
{ } {
}
}
{}{}{}
ist durch

\mathdisp {\left(  a_{ij} \cdot b_{k \ell}  \right)_{1 \leq i \leq m,\,     1 \leq  k \leq p;\,   1 \leq j \leq n,     \, 1 \leq \ell \leq r   }} {  }

gegeben.







\inputaufgabegibtloesung
{}
{





Berechne das
\definitionsverweis {Kroneckerprodukt}{}{}
der beiden Matrizen
\mathkor {} {\begin{pmatrix}  3   &   -4  \\  5  &  -2 \end{pmatrix}} {und} {\begin{pmatrix}  -2  &   7   \\  6  &  3 \end{pmatrix}} {.}





}
{} {}

\inputaufgabe
{}
{





Es sei $K$ ein
\definitionsverweis {Körper}{}{}
und

\mavergleichskettedisp
{\vergleichskette
{ A
}
{ =} { \left(  a_{ij}  \right)_{1 \leq i \leq m,\, 1 \leq j \leq n}
}
{ } {
}
{ } {
}
{ } {
}
}
{}{}{}
und

\mavergleichskettedisp
{\vergleichskette
{ B
}
{ =} { \left(  b_{k \ell }  \right)_{1 \leq  k \leq p,\, 1 \leq \ell \leq r}
}
{ } {
}
{ } {
}
{ } {
}
}
{}{}{}
\definitionsverweis {Matrizen}{}{}
mit den zugehörigen
\definitionsverweis {linearen Abbildungen}{}{}
\maabb {A} {K^n  } { K^m
} {}
bzw.
\maabb {B} {K^r } { K^p
} {.}
Zeige, dass das
\definitionsverweis {Tensorprodukt}{}{}
dieser linearen Abbildungen bezüglich der Basen

\mathbed {e_j \otimes e_\ell} {}
 {1 \leq j \leq n,\, 1 \leq \ell \leq r} {}
 {} {} {} {,}
von
\mathl{K^n \otimes K^r}{} und

\mathbed {e_i \otimes e_k} {}
 {1 \leq i \leq m,\, 1 \leq  k \leq p} {}
 {} {} {} {,}
von
\mathl{K^m \otimes K^p}{} durch das
\definitionsverweis {Kroneckerprodukt}{}{}
von $A$ und $B$ beschrieben wird.





}
{} {}

\inputaufgabe
{}
{





Zeige, dass das
\definitionsverweis {Tensorprodukt}{}{}
des
\definitionsverweis {Möbiusbandes}{}{}
mit sich selbst ein triviales Geradenbündel ist.





}
{} {}

\inputaufgabe
{}
{





Es seien
\mathkor {} {{ \mathcal  F   }} {und} {{ \mathcal  G   }} {}
\definitionsverweis {Prägarben}{}{}
auf dem
\definitionsverweis {topologischen Raum}{}{}
$X$. Zeige, dass durch
\mathl{U \mapsto { \mathcal  F   }  \left(   U   \right) \times { \mathcal  G   }  \left(   U   \right)}{} mit den natürlichen
\definitionsverweis {Produktabbildungen}{}{}
eine Prägarbe auf $X$ gegeben ist.





}
{} {}

\inputaufgabe
{}
{





Es sei $I$ eine Indexmenge und sei

\mathbed {{ \mathcal  F   }_i} {}
 {i \in I} {}
 {} {} {} {,}
eine Familie von
\definitionsverweis {Prägarben}{}{}
auf dem
\definitionsverweis {topologischen Raum}{}{}
$X$. Zeige, dass durch
\mathl{U \mapsto \prod_{i \in I} { \mathcal F_i   }  \left(   U   \right)}{} mit den natürlichen
\definitionsverweis {Produktabbildungen}{}{}
eine Prägarbe auf $X$ gegeben ist.





}
{} {}

\inputaufgabe
{}
{





Interpretiere

Beispiel 3.8

im Sinne von

Beispiel 3.12
.





}
{} {}

\inputaufgabe
{}
{





Es sei
\maabb {p} { V } { X
} {}
ein
\definitionsverweis {reelles Vektorbündel}{}{}
vom Rang $m$ auf einem
\definitionsverweis {topologischen Raum}{}{}
$X$. Zeige, dass zu jeder offenen Menge

\mavergleichskette
{\vergleichskette
{U
}
{ \subseteq }{ X
}
{  }{
}
{    }{
}
{   }{
}
}
{}{}{,}
auf der $V$ trivial ist, die zugehörige
\definitionsverweis {Prägarbe}{}{}
der
\definitionsverweis {stetigen Schnitte}{}{}
isomorph
\zusatzklammer {in welchem Sinn} {?} {}
zu
\mathl{C^0(U,\R)^m}{} ist.





}
{} {}

\inputaufgabe
{}
{





Zeige, dass die
\definitionsverweis {Gruppen}{}{}

\mathl{(\R,+)}{,}
\mathl{(\R \setminus \{0\}, \cdot)}{,}
\mathl{( {\mathbb C} ,+)}{,}
\mathl{( {\mathbb C}  \setminus \{0\}, \cdot)}{,}
\mathl{(\R^n,+)}{,}
\mathl{(S^1,\text{ mit der Winkeladdition} )}{,} die
\definitionsverweis {allgemeine lineare Gruppe}{}{}

\mathl{\operatorname{GL}_{  n   } \! \left(  \R  \right)}{} bzw.
\mathl{\operatorname{GL}_{  n   } \! \left(   {\mathbb C}   \right)}{}
\definitionsverweis {topologische Gruppen}{}{}
sind.





}
{} {}

\inputaufgabe
{}
{





Es sei $G$ eine
\definitionsverweis {topologische Gruppe}{}{}
und

\mavergleichskette
{\vergleichskette
{H
}
{ \subseteq }{ G
}
{  }{
}
{    }{
}
{   }{
}
}
{}{}{}
eine
\definitionsverweis {Untergruppe}{}{.}
Zeige, dass auf jedem
\definitionsverweis {topologischen Raum}{}{}
$X$ die
\definitionsverweis {Prägarbe}{}{}
$C^0(-,H)$ eine
\definitionsverweis {Unterprägarbe}{}{}
von $C^0(-,G)$ ist.





}
{} {}







Eine
\definitionsverweis {differenzierbare Mannigfaltigkeit}{}{}
$G$, die zugleich eine
\definitionsverweis {Gruppe}{}{}
ist, für die die Inversenabbildung und die Gruppenverknüpfung
\definitionsverweis {differenzierbare Abbildungen}{}{}
sind, heißt
\zusatzklammer {reelle} {} {}
\definitionswort {Lie-Gruppe}{.}







\inputaufgabe
{}
{





Zeige, dass die
\definitionsverweis {Gruppen}{}{}

\mathl{(\R,+)}{,}
\mathl{(\R \setminus \{0\}, \cdot)}{,}
\mathl{( {\mathbb C} ,+)}{,}
\mathl{( {\mathbb C}  \setminus \{0\}, \cdot)}{,}
\mathl{(\R^n,+)}{,}
\mathl{(S^1,\text{ mit der Winkeladdition} )}{,} die
\definitionsverweis {allgemeine lineare Gruppe}{}{}

\mathl{\operatorname{GL}_{  n   } \! \left(  \R  \right)}{} bzw.
\mathl{\operatorname{GL}_{  n   } \! \left(   {\mathbb C}   \right)}{}
\definitionsverweis {Lie-Gruppen}{}{}
sind.





}
{} {}

\inputaufgabe
{}
{





Zeige, dass das
\definitionsverweis {Tangentialbündel}{}{}
auf einer
\definitionsverweis {Lie-Gruppe}{}{}
\definitionsverweis {trivial}{}{}
ist.





}
{} {Tipp: Zeige, dass man den Tangentialraum am neutralen Element in sinnvoller Weise in die anderen Tangentialräume transportieren kann.}

\inputaufgabe
{}
{





Es sei $I$ eine
\definitionsverweis {gerichtete Indexmenge}{}{}
und sei $M_i$,
\mathl{i \in I}{,} ein
\definitionsverweis {gerichtetes System}{}{}
von Mengen. Es sei $N$ eine weitere Menge und zu jedem

\mavergleichskette
{\vergleichskette
{ i
}
{ \in }{ I
}
{  }{
}
{    }{
}
{   }{
}
}
{}{}{}
sei eine Abbildung
\maabbdisp {\psi_i} { M_i } { N
} {}
mit der Eigenschaft gegeben, dass

\mavergleichskette
{\vergleichskette
{ \psi_i
}
{ = }{ \psi_j \circ \varphi_{ij}
}
{  }{
}
{    }{
}
{   }{
}
}
{}{}{}
ist für alle

\mavergleichskette
{\vergleichskette
{ i
}
{ \preccurlyeq }{ j
}
{  }{
}
{    }{
}
{   }{
}
}
{}{}{}
\zusatzklammer {wobei $\varphi_{ij}$ die Abbildungen des Systems bezeichnen} {} {.}
Beweise die universelle Eigenschaft des
\definitionsverweis {Kolimes}{}{,}
nämlich, dass es eine eindeutig bestimmte Abbildung
\maabbdisp {\psi} { \operatorname{colim}\,_{i \in I} M_i  } { N
} {}
derart gibt, dass

\mavergleichskette
{\vergleichskette
{ \psi_i
}
{ = }{ \psi  \circ j_i
}
{  }{
}
{    }{
}
{   }{
}
}
{}{}{}
ist, wobei
\mathl{j_i:M_i \rightarrow \operatorname{colim}\,_{i \in I} M_i}{} die natürlichen Abbildungen sind.

Zeige ferner, dass falls $M_i$ ein gerichtetes System von Gruppen und falls $N$ ebenfalls eine Gruppe ist und alle $\psi_i$ Gruppenhomomorphismen sind, dass dann auch $\psi$ ein Gruppenhomomorphismus ist.





}
{} {}

\inputaufgabe
{}
{





Es sei $I$ eine
\definitionsverweis {gerichtete Indexmenge}{}{}
und sei $G_i$,
\mathl{i \in I}{,} ein
\definitionsverweis {gerichtetes System}{}{}
von kommutativen Gruppen. Zeige, dass der
\definitionsverweis {Kolimes}{}{}
eine kommutative Gruppe ist.





}
{} {}

\inputaufgabe
{}
{





Es sei $R$ ein
\definitionsverweis {kommutativer Ring}{}{}
und sei

\mavergleichskette
{\vergleichskette
{S
}
{ \subseteq }{ R
}
{  }{
}
{    }{
}
{   }{
}
}
{}{}{}
ein
\definitionsverweis {multiplikatives System}{}{.}
Auf $S$ betrachten wir folgende
\zusatzklammer {partielle} {} {}
\definitionsverweis {Ordnung}{}{,}
und zwar sagen wir

\mavergleichskette
{\vergleichskette
{f
}
{ \preccurlyeq }{g
}
{  }{
}
{    }{
}
{   }{
}
}
{}{}{,}
falls $f$ eine Potenz von $g$ teilt
\zusatzklammer {und wir identifizieren zwei Elemente, wenn diese Relation in beide Richtungen vorliegt} {} {.}
Zeige, dass die kommutativen Ringe

\mathdisp {R_f, \, f \in S} { , }

ein
\definitionsverweis {gerichtetes System}{}{}
bilden, und dass für den
\definitionsverweis {Kolimes}{}{}

\mavergleichskettedisp
{\vergleichskette
{ \operatorname{colim}_{f \in S} R_f
}
{ =} { R_S
}
{ } {
}
{ } {
}
{ } {
}
}
{}{}{}
gilt.





}
{} {}

\inputaufgabe
{}
{





Zeige, dass der
\definitionsverweis {Halm}{}{}
des
\definitionsverweis {Tangentialbündels}{}{}
einer
\definitionsverweis {differenzierbaren Mannigfaltigkeit}{}{}
in einem Punkt nur von der Dimension der Mannigfaltigkeit
\zusatzklammer {in diesem Punkt} {} {}
abhängt.





}
{} {}

\inputaufgabe
{}
{





Es seien
\mathkor {} {{ \mathcal  F   }} {und} {{ \mathcal  G   }} {}
\definitionsverweis {Prägarben}{}{}
auf dem
\definitionsverweis {topologischen Raum}{}{}
$X$ und
\mathl{{ \mathcal  F   } \times { \mathcal  G   }}{} ihre Produktprägarbe. Zeige, dass für den
\definitionsverweis {Halm}{}{}
in jeden Punkt

\mavergleichskette
{\vergleichskette
{P
}
{ \in }{ X
}
{  }{
}
{    }{
}
{   }{
}
}
{}{}{}
die Beziehung

\mavergleichskettedisp
{\vergleichskette
{ \left(  { \mathcal  F   } \times { \mathcal  G   }   \right)_P
}
{ =} { { \mathcal  F   }_P \times { \mathcal  G   }_P
}
{ } {
}
{ } {
}
{ } {
}
}
{}{}{}
gilt.





}
{} {}

\inputaufgabe
{}
{





Es sei $X$ ein
\definitionsverweis {topologischer Raum}{}{}
und seien
\mathl{{ \mathcal  F   } , { \mathcal  G   } , { \mathcal  H   }}{}
\definitionsverweis {Prägarben}{}{}
auf $X$. Zeige die folgenden Aussagen.
\aufzaehlungdrei{Die Identität
\maabbdisp {} { { \mathcal  F   } } { { \mathcal  F   }
} {}
ist ein
\definitionsverweis {Homomorphismus von Prägarben}{}{.}
}{Wenn
\maabb {\varphi} { { \mathcal F  } } { { \mathcal G  }
} {}
und
\maabb {\psi} { { \mathcal G  } } { { \mathcal H  }
} {}
Homomorphismen von Prägarben sind, so ist auch die Verknüpfung
\mathl{\psi \circ \varphi}{} ein Homomorphismus von Prägarben.
}{Zu einer
\definitionsverweis {Unterprägarbe}{}{}

\mavergleichskette
{\vergleichskette
{ { \mathcal F  }
}
{ \subseteq }{ { \mathcal G  }
}
{  }{
}
{    }{
}
{   }{
}
}
{}{}{}
ist die natürliche Inklusion ein Homomorphismus von Prägraben.
}





}
{} {}

\inputaufgabe
{}
{





Es sei $I$ eine Indexmenge und sei

\mathbed {{ \mathcal  F   }_i} {}
 {i \in I} {}
 {} {} {} {,}
eine Familie von
\definitionsverweis {Prägarben}{}{}
auf dem
\definitionsverweis {topologischen Raum}{}{}
$X$ mit der Produktprägarbe
\mathl{\prod_{i \in I} { \mathcal F_i   }}{.} Es sei ${ \mathcal  G   }$ eine weitere Prägarbe auf $X$. Zeige, dass ein
\definitionsverweis {Prägarbenmorphismus}{}{}
\maabbdisp {\psi} { { \mathcal  G   } } { \prod_{i \in I} { \mathcal F_i   }
} {}
das gleiche ist wie eine Familie von Prägarbenmorphismen
\maabbdisp {\psi_i} { { \mathcal  G   } } { { \mathcal F_i   }
} {.}





}
{} {}
